% ============================================================================
%  Replacement for Section 6 -- Cross-Dataset Validation on CBIS-DDSM
%
%  PREAMBLE REQUIREMENTS (add to your main .tex if not already present):
%     \usepackage{booktabs}
%     \usepackage{amssymb}
%     \usepackage{pifont}
%     \newcommand{\cmark}{\ding{51}}
%     \newcommand{\xmark}{\ding{55}}
%
%  All figures traceable to W&B project maya-hayat-ariel-university/ddsm-acgan.
%  Supporting detail: DDSM-ACGAN-Pytorch/DDSM_RESULTS.md
% ============================================================================

\section{Cross-Dataset Validation on CBIS-DDSM}
\label{sec:ddsm}

\subsection{Motivation and Task}

A central concern in any empirical machine learning study is data specificity. Is our failure to
reproduce the paper's accuracy increase an artifact of the modern, cleaner COVID-CXR dataset, or is it
a property of the AC-GAN method itself under frozen-backbone constraints? To answer this, we
replicated the entire pipeline on \textbf{CBIS-DDSM}, a static and versioned breast mammography
benchmark. We restricted attention to mass-type ROI lesions, giving \textbf{1{,}318 real training
images and 378 test images} (231 benign, 147 malignant). The always-predict-benign floor is
\textbf{61.11\%}; no result below this threshold carries information.

CBIS-DDSM is a useful contrast to COVID-CXR precisely because it is a \emph{hard} task. Our COVID
baseline sits near ceiling at 94.48\%, leaving little headroom for augmentation to fill. DDSM leaves
a great deal.

\subsection{Experimental Setup}

We ran the same three-component pipeline (AC-GAN generator, discriminator, and VGG16 classifier) with
identical hyperparameters, and applied both Stage 2 tracks: the generator improvements (spectral
normalization on every discriminator convolution; transpose-convolution kernels changed from 5 to 4 at
stride 2) and the classifier capacity change (unfreezing the top two VGG16 blocks, a BatchNorm head,
and discriminative learning rates).

\paragraph{Data adjustment.} DDSM ROIs are stored as DICOM and converted to 16-bit grayscale PNGs with
pixel values up to 65{,}535. Standard PIL conversion truncates to 8 bits, collapsing nearly every pixel
to flat white. We implemented per-image min--max scaling to $[0,255]$ prior to RGB conversion,
restoring tissue texture and rescuing the baseline classifier. All results below post-date this fix.

\paragraph{A measurement problem that had to be solved first.} Our initial DDSM experiments were run
without fixed random seeds. When we introduced seeding and re-ran them, several results changed sign.
Investigating this, we isolated three independent sources of run-to-run variance, none of which had
been controlled (Table~\ref{tab:ddsm_variance}).

\begin{table}[h]
\centering
\caption{Sources of run-to-run variance on CBIS-DDSM. The synthetic pool draw dominates, and exceeds
every effect size we were attempting to measure.}
\label{tab:ddsm_variance}
\begin{tabular}{llr}
\toprule
\textbf{Source} & \textbf{Evidence} & \textbf{Magnitude} \\
\midrule
Classifier random seed & AD unfrozen, seeds 0/1/2: 69.84 / 69.05 / 70.90 & 1.85 pts \\
\textbf{Synthetic pool noise draw} & Same checkpoint, different $z$ sample & \textbf{5.29 pts} \\
Pool regenerated on another GPU & Identical seed and code, different session & 2.65 pts \\
\bottomrule
\end{tabular}
\end{table}

Since \textbf{no effect we sought to measure exceeded five points}, any single-run comparison on this
dataset is dominated by sampling noise. This invalidated several of our own earlier DDSM comparisons,
which we withdraw.

The third source merits comment for reproducibility practice generally. Classifier training is fully
deterministic given a seed: our real-data-only runs reproduce to the digit across sessions five days
apart. Synthetic \emph{generation}, however, is not, because cuDNN selects convolution algorithms per
device. \textbf{A seeded synthetic pool is reproducible only on the machine that created it.}

Accordingly, every figure reported below is the mean of \textbf{three independent synthetic pool
draws} with the classifier seed held fixed, or the mean of three classifier seeds where no synthetic
pool is involved.

\subsection{Main Results}

Synthetic pools were drawn from each architecture's best checkpoint as measured by KID (baseline epoch
650, improved epoch 750), with 300 malignant and 300 benign images per pool. Results are given in
Table~\ref{tab:ddsm_main}. We report \emph{balanced accuracy} --- the mean of the two class recalls ---
because raw accuracy is misleading on an imbalanced test set, and because it is precisely what
distinguishes the two findings that follow.

\begin{table}[h]
\centering
\caption{CBIS-DDSM main results. Each row averages three independent runs. Balanced accuracy is the
mean of malignant and benign recall.}
\label{tab:ddsm_main}
\begin{tabular}{llrrrrr}
\toprule
\textbf{Mode} & \textbf{Classifier} & \textbf{Train acc} & \textbf{Test acc} &
\textbf{Mal.\ recall} & \textbf{Ben.\ recall} & \textbf{Bal.\ acc} \\
\midrule
SA (improved GAN) & frozen & 0.820 & 62.70 & 0.610 & 0.637 & 0.624 \\
SA (baseline GAN) & frozen & 0.848 & 63.80 & 0.556 & 0.691 & 0.623 \\
\textbf{AD (real only)} & frozen & 0.795 & \textbf{66.05} & 0.535 & 0.740 & \textbf{0.638} \\
\midrule
SA (improved GAN) & uf${=}2$ & 0.997 & 67.55 & 0.712 & 0.652 & 0.682 \\
SA (baseline GAN) & uf${=}2$ & 0.997 & 68.21 & 0.706 & 0.667 & 0.686 \\
\textbf{AD (real only)} & uf${=}2$ & 0.993 & \textbf{69.93} & 0.646 & 0.733 & \textbf{0.690} \\
\bottomrule
\end{tabular}
\end{table}

\subsection{Finding 1: Unfreezing the Encoder Produces Genuine Learning}

Comparing the two real-data-only rows of Table~\ref{tab:ddsm_main} isolates the effect of classifier
capacity:

\begin{table}[h]
\centering
\caption{Effect of unfreezing the encoder on real-data-only training. Both class recalls and balanced
accuracy improve together.}
\label{tab:ddsm_unfreeze}
\begin{tabular}{lrrr}
\toprule
\textbf{Metric} & \textbf{AD frozen} & \textbf{AD uf${=}2$} & \textbf{Change} \\
\midrule
Test accuracy      & 66.05 & 69.93 & $+3.88$ \\
Malignant recall   & 0.535 & 0.646 & $\mathbf{+0.111}$ \\
Benign recall      & 0.740 & 0.733 & $-0.007$ \\
Balanced accuracy  & 0.638 & 0.690 & $\mathbf{+0.052}$ \\
\bottomrule
\end{tabular}
\end{table}

Unfreezing correctly identifies \textbf{16 additional malignant cases out of 147} while losing
essentially no benign cases. Both class recalls and balanced accuracy improve together, which is the
signature of a model that has genuinely learned better features rather than one that has merely
relocated its decision threshold. This confirms the \emph{Capacity Knob} of
Section~\ref{sec:discussion} independently on a second medical domain: ImageNet features must be
permitted to adapt to medical texture.

\subsection{Finding 2: Augmentation Shifts the Operating Point Without Improving Discrimination}

Every augmented configuration scores \emph{below} the real-data-only baseline of matched capacity
(Table~\ref{tab:ddsm_deltas}).

\begin{table}[h]
\centering
\caption{Change induced by synthetic augmentation, relative to the real-data-only baseline of the
same classifier capacity. Malignant-recall gains are offset almost exactly by benign-recall losses.}
\label{tab:ddsm_deltas}
\begin{tabular}{llrrrr}
\toprule
\textbf{Classifier} & \textbf{Synthetic pool} & $\Delta$\textbf{Acc} &
$\Delta$\textbf{Mal.\ rec} & $\Delta$\textbf{Ben.\ rec} & $\Delta$\textbf{Bal.\ acc} \\
\midrule
frozen   & baseline GAN & $-2.25$ & $+0.020$ & $-0.049$ & $-0.015$ \\
frozen   & improved GAN & $-3.35$ & $+0.075$ & $-0.103$ & $-0.014$ \\
uf${=}2$ & baseline GAN & $-1.72$ & $+0.060$ & $-0.067$ & $-0.003$ \\
uf${=}2$ & improved GAN & $-2.38$ & $+0.066$ & $-0.081$ & $-0.008$ \\
\bottomrule
\end{tabular}
\end{table}

Malignant recall does rise consistently, by as much as $+0.075$, and on a cancer-detection task this
is the recall of clinical consequence. It is tempting to read this as evidence that the model has
learned malignant features more effectively. \textbf{The balanced-accuracy column shows that it has
not.} In every configuration the gain in malignant recall is offset almost exactly by a loss in benign
recall, leaving balanced accuracy flat or marginally negative. In absolute case counts, the unfrozen
baseline-GAN configuration identifies \textbf{9 additional malignant cases} while misclassifying
\textbf{14 additional benign cases}.

This is a decision-threshold shift rather than improved discrimination, and it has a mundane
explanation: adding 300 synthetic malignant images to a training set with a 231:147 benign majority
moves the class prior. It requires nothing to have been learned from the synthetic images. The
contrast with Finding 1 is instructive --- unfreezing raised malignant recall \emph{at no cost},
whereas augmentation merely traded for it.

We note that this trade-off may nevertheless be operationally desirable. In a screening context a
missed malignancy carries greater cost than a false positive, so a more sensitive operating point can
be preferred even at equal balanced accuracy. That, however, is a statement about which point on the
trade-off curve one wishes to occupy, not evidence that augmentation improved the underlying model ---
and the same shift is obtainable at no cost through class weighting or threshold adjustment, with no
generator required.

\subsection{Finding 3: The Generator Improvements Trade Quality Ceiling for Stability}

We trained both generator architectures to 900 epochs and evaluated \emph{every} saved checkpoint
(18 per architecture) using KID against the real training set.

\paragraph{Training stability improved unambiguously}, reproducing the Track A result from the COVID
experiments:

\begin{table}[h]
\centering
\caption{Loss variability by architecture and training window (standard deviation of per-epoch loss).}
\label{tab:ddsm_stability}
\begin{tabular}{lrr}
\toprule
\textbf{Window} & \textbf{$D_{\text{loss}}$ sd} & \textbf{$G_{\text{loss}}$ sd} \\
\midrule
Baseline, full 300 epochs   & 0.351 & 0.984 \\
Improved, epochs 201--300   & \textbf{0.031} & \textbf{0.025} \\
Improved, epochs 351--600   & \textbf{0.026} & \textbf{0.018} \\
\bottomrule
\end{tabular}
\end{table}

\paragraph{Generation quality, however, did not improve.} Table~\ref{tab:ddsm_kid} reports KID across
training for both architectures.

\begin{table}[h]
\centering
\caption{KID across training (lower is better). The baseline reaches a superior ceiling at epoch 650
but diverges catastrophically thereafter; the improved architecture never diverges but never matches
that ceiling.}
\label{tab:ddsm_kid}
\begin{tabular}{lrrrrrrrr}
\toprule
\textbf{Epoch} & 250 & 400 & 500 & 600 & 650 & 700 & 750 & 900 \\
\midrule
Baseline & 0.164 & 0.170 & 0.158 & 0.134 & \textbf{0.121} & 0.419 & 0.437 & 0.404 \\
Improved & 0.390 & 0.283 & 0.220 & 0.203 & 0.259 & 0.206 & \textbf{0.194} & 0.264 \\
\bottomrule
\end{tabular}
\end{table}

Two opposing facts follow. First, \textbf{the baseline architecture reaches a higher quality ceiling}:
its best checkpoint (KID 0.121) surpasses the improved architecture's best (0.194) using 350 fewer
epochs, and it wins 12 of 18 matched-epoch comparisons. Second, \textbf{the baseline architecture is
not safe to train unattended}: at epoch 684 it diverged catastrophically, with KID rising 245\% within
50 epochs and never recovering. The loss trace identifies the mechanism precisely --- $D_{\text{loss}}$
fell from 1.17 to 0.38 while $G_{\text{loss}}$ rose from 2.32 to 7.81, indicating that the
discriminator overpowered the generator until the generator's gradients vanished. The improved
architecture held $D \approx 1.60$ and $G \approx 0.90$ flat from epoch 200 through 900 and never
diverged.

This is precisely the mechanism spectral normalization implements. Bounding the discriminator's
Lipschitz constant \emph{weakens} $D$ by construction: a weaker discriminator conveys less information
to the generator, costing peak quality, but it can never overpower the generator.
\textbf{Lower ceiling, no floor.} Evaluated against a pre-registered criterion (median KID over the
final six checkpoints) the improved architecture wins, 0.248 to 0.411; evaluated on best-ever
checkpoint the baseline wins, 0.121 to 0.194. Both statements are true, and they answer different
questions: if one monitors quality and stops early, the baseline is preferable; if one trains to a
fixed budget and accepts the result, the improved architecture is preferable.

Methodologically, adjacent checkpoints 50 epochs apart differ by as much as 245\% in KID against error
bars of $\pm 0.015$. Comparing a single checkpoint per architecture --- as we initially did --- measures
where training happened to be sampled rather than the architecture itself.

\subsection{Finding 4: Generation Quality Is Decoupled from Augmentation Value}

CBIS-DDSM permits the cleanest available test of the decoupling claim, because it admits a comparison
with \emph{no architecture confound}: the same generator at two training lengths
(Table~\ref{tab:ddsm_decoupling}).

\begin{table}[h]
\centering
\caption{Decoupling of generation quality from augmentation value, with no architecture confound.}
\label{tab:ddsm_decoupling}
\begin{tabular}{lrr}
\toprule
\textbf{Baseline generator} & \textbf{KID} & \textbf{Unfrozen CNN-SA accuracy} \\
\midrule
Epoch 300 & 0.315 & 68.96 \\
Epoch 650 & \textbf{0.121} & 68.61 \\
\bottomrule
\end{tabular}
\end{table}

\textbf{A $2.6\times$ improvement in generation quality changed downstream accuracy by $-0.35$ points
--- that is, not at all.} This is the DDSM counterpart to the COVID Stage 1 result in which FID halved
with no downstream movement, and it constitutes stronger evidence because only a single variable
changed.

Consistently, no DDSM generator produced transferable pathology. Training solely on a synthetic pool
and evaluating on real data never meaningfully cleared the 61.11\% majority floor (60.05\%, 59.79\%,
and 65.08\% for the baseline, improved-300 and improved-600 pools respectively), and a
real-versus-synthetic discriminability probe separated \emph{every} pool from real data with 100\%
held-out accuracy.

We are careful not to overclaim. On DDSM the rankings of generator quality and augmentation value
happen to \emph{agree} across architectures: the improved generator is worse on KID and worse
downstream. We therefore \textbf{do not} claim that DDSM reproduces the COVID Reversal Paradox, in
which a strictly superior generator degraded the classifier; on DDSM no generator was strictly
superior. What DDSM establishes is the weaker but still decisive claim that \textbf{generator quality
metrics carry no information about augmentation value in either direction}.

\subsection{Cross-Dataset Summary}

\begin{table}[h]
\centering
\caption{Cross-dataset validation matrix. DDSM replicates four findings, refines one, and contradicts
two.}
\label{tab:cross_dataset_validation}
\begin{tabular}{lcc}
\toprule
\textbf{Finding} & \textbf{COVID-CXR} & \textbf{CBIS-DDSM} \\
\midrule
Naive paper setup fails to lift accuracy & \cmark & \cmark \\
Unfreezing the encoder raises the baseline & \cmark & \cmark\ ($+3.88$) \\
Naive synthetic pool carries no transferable pathology & \cmark & \cmark \\
\textbf{Unfreezing rescues the augmentation lift} & \cmark\ ($+2.19$) & \xmark\ $\mathbf{(-1.72)}$ \\
Generator quality decoupled from augmentation value & \cmark & \cmark\ (strongest evidence) \\
Improved architecture stabilises GAN training & \cmark & \cmark \\
Improved architecture raises generation quality & \cmark\ (FID $273 \to 121$) & \xmark \\
\bottomrule
\end{tabular}
\end{table}

DDSM replicates four of our seven findings, refines one, and contradicts two. The contradictions are
the informative part: augmentation's success on COVID-CXR is not a general property of the method, and
the generator modifications that improved quality on chest X-rays purchased stability rather than
quality on mammography.


% ============================================================================
%  REPLACEMENT TEXT FOR OTHER SECTIONS
%  (copy the relevant paragraph into the existing section; not meant to be
%   compiled as-is in place)
% ============================================================================

% ---------------------------------------------------------------- ABSTRACT --
% Replace the final paragraph's DDSM sentences with:
%
% The results reproduced our core diagnostic findings: the naive paper setup
% failed to improve accuracy, unfreezing the encoder was again the only
% intervention that reliably improved the model (+3.88%, with malignant recall
% rising 0.111 at no cost to benign recall), and generator quality was again
% decoupled from downstream value, with a 2.6x KID improvement producing no
% accuracy change. Critically, DDSM also bounds our central claim: augmentation
% did not help under any configuration tested, including one built from the
% best generator we produced. Where augmentation did shift malignant recall
% upward, balanced accuracy remained flat -- a decision-threshold shift rather
% than improved discrimination. Unfreezing is therefore necessary but not
% sufficient. Establishing this also required quantifying a measurement problem
% that invalidated several of our own earlier DDSM comparisons: uncontrolled
% synthetic-pool sampling contributes up to 5 accuracy points, exceeding every
% effect we sought to measure.

% ------------------------------------------------------------ 7.1 WHAT WORKED
% Scope the claim that capacity expansion "rescued the paper's augmentation
% lift" to the COVID dataset. On DDSM capacity raised the baseline but produced
% no lift.

% ------------------------------------------------------- 7.2 WHAT DID NOT WORK
% Add:
%
% On CBIS-DDSM, augmentation failed to produce a lift under every configuration
% tested, including an unfrozen classifier paired with our best-quality
% generator. Where malignant recall improved, balanced accuracy did not: the
% gain was a threshold shift obtainable by class weighting alone. Ample headroom
% was not sufficient to rescue augmentation.

% ------------------------------------------------- 7.3 DUAL-KNOB FRAMEWORK ---
% Add a boundary condition to the Headroom Knob:
%
% Headroom is necessary but not sufficient. CBIS-DDSM offered abundant headroom
% (69.93% against a 61.11% floor) and augmentation still failed, because the
% generator never produced class-discriminative samples: trained alone, its
% pools could not beat the majority-class floor. A third precondition --
% generator transfer quality above the majority floor -- must hold before
% capacity and headroom can matter. This also supplies a cheap go/no-go test:
% run the transfer probe first; if it does not clear the floor, no downstream
% configuration will rescue the augmentation.
%
% And to the decoupling paragraph:
%
% This was confirmed on DDSM under the cleanest available control -- the same
% generator architecture at two training lengths, where a 2.6x KID improvement
% produced a -0.35 point accuracy change.
